\subsection{22. Generative AI-Driven Accelerated Discovery of Passivation Molecules for Perovskite Solar Cells}
\textbf{OSTI:} \texttt{PVMol-Gen-Fajar2026}\quad\textbf{Authors:} Fajar et al. (2026)\quad\textbf{Domain:} Materials/Generative ML (perovskite passivation)\\
\scorebadge{Coverage}{7}\quad\scorebadge{Agreement}{5}\quad\textbf{Total:} 12/20\par\smallskip
\noindent\textbf{Coverage rationale.} Full three-stage pipeline implemented (SMILES-X classifier, iterative GPT-2 generation for 3 cycles at paper scale of 100K/cycle, 7 RDKit filters + clustering). Ran on real 8$\times$A100 hardware. Missing: xTB/DFT E\_gap and dipole filtering (used RDKit surrogates), and paper's exact SMILES tokenization pathway (we used SELFIES variant in production run).\par
\noindent\textbf{Agreement rationale.} Generation rates reasonable (83-88\% class-1 per cycle). Key quantitative gap: Stage-1 classifier F1=0.656$\pm$0.031 vs paper 0.80, ROC-AUC 0.620 vs 0.88 --- major reproducibility concern pointing to library/hyperparameter differences we could not pin down. SELFIES pathway yields 6.6$\times$ more filtered candidates than SMILES, so downstream numbers diverge by design. Final 10-cluster representatives produced but not externally validated.\par\smallskip
\noindent\textbf{What reproduced:}
\begin{itemize}[leftmargin=1.4em,itemsep=0.2em,topsep=0.2em]
\item SMILES-X-style discriminative classifier with 5-fold CV on 314 labeled molecules
\item PubChem similarity expansion ($\geq$80\% Tanimoto) to build \textasciitilde{}11K training set
\item GPT-2 fine-tuning + generate-classify-retrain loop for 3 cycles at 100K molecules/cycle
\item Stage-3 RDKit filters (SA score, PAINS, HBD/HBA, TPSA)
\item Clustering into 10 groups with representative selection
\item Overall pipeline architecture end-to-end at paper scale
\end{itemize}\noindent\textbf{What missing:}
\begin{itemize}[leftmargin=1.4em,itemsep=0.2em,topsep=0.2em]
\item Matching classifier metrics (F1 0.656 vs 0.80) --- root cause not identified
\item xTB/DFT energy gap and dipole filters (used RDKit approximations)
\item Paper's exact SMILES augmentation protocol (we swapped to SELFIES mid-pipeline)
\item External experimental validation of top candidates
\item Synthetic accessibility tie-breaking rules applied by paper for final down-select
\item DFT-level verification of E\_gap$\in$[1.5,5.0] eV and dipole$\in$[1.5,4.0] D per molecule
\end{itemize}\noindent\textbf{Follow-on research questions:}
\begin{enumerate}[leftmargin=1.6em,itemsep=0.2em,topsep=0.2em,label=Q\arabic*.]
\item Does matching the paper's exact SMILES-X version + tokenizer close the F1 gap from 0.656 to 0.80, or is the gap due to training-data preprocessing (canonicalization, augmentation multiplier)
\item Run xTB on the 10 final candidates from both SMILES and SELFIES pipelines --- do the E\_gap/dipole distributions overlap with the paper's reported 10 molecules
\item Does training-set contamination from predicted class-1 relabeling (cycles 2-3) inflate apparent classifier accuracy Quantify by holding out a fixed external test set and tracking its F1 across cycles.
\item Substitute the GPT-2 generator with a chemistry-tuned foundation model (e.g., ChemGPT, MolGPT) --- does the class-1 rate at cycle 1 jump from 83\% to \textgreater{}95\%, and do top-10 clusters converge to similar motifs
\item Benchmark downstream perovskite figure-of-merit predictors (DFT binding energy on Cs/FA surfaces) on the 20 candidates (10 SMILES + 10 SELFIES) to test whether the representation choice biases toward different chemistries.
\end{enumerate}

